\documentclass[lang=cn,newtx,10pt,scheme=chinese]{elegantbook}

\title{C 语言课程讲义}

\author{LMXZ}
\date{2025-10-21}
\version{0.1}

\setcounter{tocdepth}{3}

\cover{cover.jpg}

% 本文档命令
\usepackage{array}
\usepackage{titlesec}

\newcommand{\codex}{\lstinline}

% 修改标题页的橙色带
\definecolor{customcolor}{RGB}{32,178,170}
\colorlet{coverlinecolor}{customcolor}
\usepackage{cprotect}

\addbibresource[location=local]{reference.bib} % 参考文献，不要删除

\begin{document}

\maketitle
\frontmatter

\tableofcontents

\mainmatter

\chapter{快速入门}

\section{快速配置开发工具}

前期建议使用易配置的 Code::Blocks 作为编程工具。可以在这里下载：

\href{https://www.codeblocks.org/downloads/binaries/}
{https://www.codeblocks.org/downloads/binaries/}

选择“Windows(64 bit)”（如果你用的是 Windows 的话）下的 \lstinline{codeblocks-X.Xmingw-setup.exe}。这个版本自带了 GCC 编译器，安装后即可使用。

使用 Code::Blocks 只是为了快速上手，但并不推荐使用作为长期的开发工具，等你熟悉了 C 语言之后，可以考虑使用更强大的 IDE，比如 Visual Studio Code、CLion 等。

\section{第一个程序，Hello World！}

Hello World 是指输出“Hello, World”的简单程序，在开始学习一门语言时，一般会从这里开始入手。

\begin{lstlisting}[language=C]
    #include <stdio.h>
    
    int main() {
        printf("Hello, World!"); // 输出 Hello World
        return 0; // 程序结束
    }
\end{lstlisting}        

现在程序写好了，需要运行时就点击上面的“\includegraphics{image/btn_compile_run.png}”按钮。

目前你需要了解的是：

\begin{itemize}
    \item \lstinline{#include <stdio.h>} 是预处理指令，用于包含标准输入输出库，提供 \lstinline{printf} 函数的定义；
    \item \lstinline{int main()} 是主函数的定义，程序从这里开始执行；
    \item \lstinline{printf("Hello, World!");} 用于输出内容到屏幕，括号内的内容是要输出的字符串；
    \item \lstinline{return 0;} 表示程序正常结束，返回值为 0；
    \item “\lstinline{// 输出 Hello World}” 这种字样是注释，用于做批注，以便人类能看懂；计算机并不会理会“\lstinline{//}”以及后面的内容；
    \item 在 \lstinline{int main} 里面的每一行的结尾都必须有一个分号 “\lstinline{;}”；
    \item 输入程序的时候要特别注意，把输入法关掉，否则你输入的符号会变成全角字符；程序中的符号应该统一为半角字符。
\end{itemize}

\section{输入输出}

刚才的例子中，\lstinline{printf("XX")} 是用于输出内容的。
然而一个程序往往需要输入、计算、输出三步。我们看第二个例子：

\label{sec:input-output}
\begin{lstlisting}[language=C]
    #include <stdio.h>
    
    int main() {
        int a, b; // 定义变量 a b；
        printf("Please input the value of a and b:");
        scanf("%d%d", &a, &b); // 读入 a b
        int c = a + b; // 计算 a+b 的结果，存到变量 c 里面
        printf("The result of %d + %d is: %d", a, b, c); // 输出结果
        return 0;
    }
\end{lstlisting}

这个程序会读取用户输入的两个数字，相加并输出。
运行程序，输入“\lstinline{3 5}”（注意中间有空格），你会看到程序输出了“\lstinline{The result of 3 + 5 is: 8}”。

这里你需要记住三句话的用法：

\begin{itemize}
    \item \lstinline{int} 用于定义整数型变量，一行可以定义多个变量；定义时如果在后面加上 “\lstinline{= xx}”，就会同时给它赋值。
    \item \lstinline{scanf} 用于输入，有两个 \lstinline{%d}，这个东西被称为“占位符”；要读取几个变量，引号中就要有多少个 \lstinline{%d}；
    \item \codex{printf} 用于输出变量，其中 \lstinline{%d} 也是占位符，每个占位符会被依次替换为后面的每个表达式的值。
\end{itemize}

\section{分支选择、数据类型、算术运算}

\begin{lstlisting}[language=C]
    #include <stdio.h>
    #include <math.h> // 引入数学库，使用 sqrt 函数

    int main() {
        double a, b, c; // 定义三个双精度浮点数变量
        printf("Solve the equation ax^2 + bx + c = 0\n");
        printf("Please input the coefficients a, b and c:");
        scanf("%lf%lf%lf", &a, &b, &c); // 读入 a b c
        double delta = b * b - 4 * a * c; // 计算判别式
        if (delta < 0) {
            printf("No real roots.\n");
        } else if (delta == 0) {
            double root = -b / (2 * a);
            printf("One real root: %lf\n", root);
        } else {
            double root1 = (-b + sqrt(delta)) / (2 * a);
            double root2 = (-b - sqrt(delta)) / (2 * a);
            printf("Two real roots: %lf and %lf\n", root1, root2);
        }
        return 0;
    }
\end{lstlisting}

这个例子展示了如何使用分支选择语句 \lstinline{if-else} 来处理不同的情况。
相比于前面的例子，这里你需要了解的内容有：

\begin{itemize}
    \item \lstinline{double} 用于定义双精度浮点数变量，可以存储小数，与 \lstinline{int} 不同；
    \item 在 \lstinline{scanf} 和 \lstinline{printf} 中，使用 \lstinline{%lf} 来表示双精度浮点数的占位符；
    \item 使用算术运算符 “\lstinline{+}”、“\lstinline{-}”、“\lstinline{*}”、“\lstinline{/}”、“\lstinline{%}” 进行加减乘除和取模（即除法求余数）运算；
    \item 使用数学库中的函数，比如这里的 \lstinline{sqrt()} 用于计算平方根；使用这些函数前需要包含相应的头文件：\lstinline{#include <math.h>}；其中还有很多常用的数学函数，如 \lstinline{pow()}（幂运算）、\lstinline{sin()}、\lstinline{cos()}（三角函数）等；
    \item 使用 \lstinline{if-else} 语句根据条件执行不同的代码块。
\end{itemize}

\section{Online Judge 练习平台}

Online Judge（OJ）是指在线编程练习平台，提供了大量的编程题目供练习。
常见的 OJ 平台有：洛谷、Codeforces、牛客网、LeetCode 等。

简单介绍一下洛谷平台的使用方法：

\begin{itemize}
    \item 访问洛谷网站：\href{https://www.luogu.com.cn/}{https://www.luogu.com.cn/}，注册并登录；
    \item 浏览题库，选择适合自己水平的题目进行练习；
    \item 阅读题目描述，理解题意和输入输出要求；
    \item 在本地编写代码，调试通过后，将代码提交到洛谷平台；
    \item 提交后，系统会自动评测代码的正确性和效率，并给出结果反馈：AC 表示正确，其他结果表示有错误或超时等问题；
    \item 选择提交时要选择 C 语言作为编程语言；
    \item 严格遵守输入输出格式，不要输出多余内容；比如 \ref{sec:input-output} 节中的程序就不适合直接提交到 OJ 平台，因为它包含了提示信息，去除“Please input...”和“The result...”，只输出一个数字表示结果即可。
\end{itemize}

\section{学习方法}

这里能讲到的内容毕竟有限，学习过程中需要善于查找资料，遇到问题时可以：
\begin{itemize}
    \item 查阅 C 语言相关书籍和在线文档，如 \href{https://en.cppreference.com/w/c}{C Reference}；这是最靠谱，最权威的参考资料；
    \item 利用搜索引擎查找相关问题的解决方案；
    \item 查阅编程社区和论坛，向其他程序员请教问题；
    \item AI 助手（比如 ChatGPT）也是一个不错的选择，可以用来解答编程相关的问题。
\end{itemize}
擅于自主学习和解决问题是成为优秀程序员的重要能力。

\begin{problemset}[练习]
    \item 在洛谷通过以下题目：
        \href{https://www.luogu.com.cn/problem/P1001}{P1001 A+B Problem}、
        \href{https://www.luogu.com.cn/problem/P1150}{P1150 Peter 的烟}、
        \href{https://www.luogu.com.cn/problem/P1425}{P1425 小鱼的游泳时间}。
    \item 现在你已经掌握了基本的输入、计算、输出三步；接下来请你编写一个程序，读取平面上三个点的坐标，计算并输出它们构成的三角形的三条边的长度、三个角的角度以及三角形面积。
\end{problemset}


\newpage

\textbf{参考答案：}

\begin{lstlisting}[language=C]
    #include <stdio.h>
    #include <math.h>

    int main() {
        const double PI = acos(-1.0);
        double x1, y1, x2, y2, x3, y3;
        printf("Please input the coordinates of three points (x1 y1 x2 y2 x3 y3):");
        scanf("%lf%lf%lf%lf%lf%lf", &x1, &y1, &x2, &y2, &x3, &y3);

        double a = sqrt((x2 - x3) * (x2 - x3) + (y2 - y3) * (y2 - y3));
        double b = sqrt((x1 - x3) * (x1 - x3) + (y1 - y3) * (y1 - y3));
        double c = sqrt((x1 - x2) * (x1 - x2) + (y1 - y2) * (y1 - y2));

        double angleA = acos((b*b + c*c - a*a) / (2*b*c)) * 180.0 / PI;
        double angleB = acos((a*a + c*c - b*b) / (2*a*c)) * 180.0 / PI;
        double angleC = 180.0 - angleA - angleB;

        double s = (a + b + c) / 2;
        double area = sqrt(s * (s - a) * (s - b) * (s - c));

        printf("Sides: a=%.2lf, b=%.2lf, c=%.2lf\n", a, b, c);
        printf("Angles: A=%.2lf°, B=%.2lf°, C=%.2lf°\n", angleA, angleB, angleC);
        printf("Area: %.2lf\n", area);

        return 0;
    }
\end{lstlisting}

\chapter{循环、数组、函数}

\section{循环}

\subsection{for循环}

\lstinline{for} 循环是最常用的循环结构，特别适合在已知循环次数的情况下使用。

\begin{lstlisting}[language=C]
    for (int i = 1; i <= 5; i++) {
        printf("%d ", i);
    }
\end{lstlisting}

语法说明：

\lstinline{int i = 1}：循环变量初始化

\lstinline{i <= 5}：循环条件

\lstinline{i++}：循环变量更新

在较老版本的 C（如 C89/ANSI C）标准下，不能在 for 语句头中声明循环变量，需要在循环外先声明变量，例如：

\begin{lstlisting}[language=C]
    int i, j;
    for (i = 1; i <= 5; i++) {
        for (j = 1; j <= 5; j++)
        printf("%d, %d\n", i, j);
    }
\end{lstlisting}

注意：这样声明的 i 在循环结束后仍然可见（作用域为所在块），如需限制作用域可另起一个内层块。

\subsection{while循环}

\lstinline{while} 循环在不确定循环次数，但知道循环条件时使用。

\begin{lstlisting}[language=C]
    int count = 1;
    // 同样打印数字1到5
    while (count <= 5) {
        printf("%d ", count);
        count++;
    }
\end{lstlisting}


\subsection{do-while循环}

\lstinline{do-while} 循环与 \lstinline{while} 循环类似，但会先执行一次循环体，再检查条件。

\begin{lstlisting}[language=C]
    int count = 1; // 把 count 改为 6 或以上的值，循环体代码也至少会执行一次
    do {
        printf("%d ", num);
        num++;
    } while (num <= 5);
\end{lstlisting}

\section{数组}

\subsection{数组声明}

在不同的 C 标准中，数组的定义方式有区别。
在 C99 及更老的版本中，数组大小只能是常量，而之后的版本中可以为变量。

\begin{lstlisting}[language=C]
    const int BUF_LEN = 16;
    int fixed_buf1[10]; // 常量长度
    int fixed_buf2[BUF_LEN]; // 符号常量也是可以的
    // 以下定义方式在老版本是不可用的
    int len1 = 16, len2 = 0;
    scanf("%d", &len2);
    int buf1[len1], buf2[len2]; // 无论是自己赋值的变量还是用户输入的变量，都可以
\end{lstlisting}

动态分配是通用的数组定义方式：

\begin{lstlisting}[language=C]
    #include <stdlib.h>

    int main {
        int n = 0;
        scanf("%d", &n);
        int *p = malloc(n * sizeof(int));
        if (p == NULL) {
            // 内存分配失败的处理
        }
        for (int i = 0; i < n; ++i)
            scanf("%d", &p[i])
        // 一顿操作以后，如果数组不再使用了，记得释放内存！！
        free(p);
        return 0;
    }
\end{lstlisting}

兼容性提示：若需要同时支持可能不含可变长度数组的实现，可以用预处理判断并退回到 malloc 方案：
\begin{lstlisting}[language=C]
    /* 如果实现不支持 VLA，则 __STDC_NO_VLA__ 会被定义 */
    #ifndef __STDC_NO_VLA__
        int arr[n];
    #else
        int *arr = malloc(n * sizeof *arr);
        /* 使用后记得 free(arr); */
    #endif
\end{lstlisting}

好复杂是不是？\textbf{现在先记住第一个示例就行了。}作为平时的练习，你在做 OJ 上的题目时一般不会使用动态分配的方式，但在一些实际工程中动态分配是到处都有的。

\subsection{数组访问}

数组下标从0开始，到 $n-1$ 结束（$n$ 为数组长度），\textbf{不包含 $n$}。
当你访问了超过这个范围的下标，就会发生越界。
如果你想要让数组下标从 $1$ 开始，可以把数组长度开得比实际使用的大小大一点。
和现代的各种语言不同，默认情况下，程序\textbf{不会}自动帮你检查数组下标是否越界并报错，
当你越界时，会出现很多匪夷所思的 bug。

那么有没有不默认的情况？有的！这一点在绝大多数国内外教材中都不会提及，但我们后面会介绍。

\section{函数}

\subsection{前置知识：表达式}

C 程序中，几乎一切操作和运算都是通过表达式实现的。举个例子：

\begin{lstlisting}[language=C]
    int a = 0;
    int b = a++;
    int c = ++a;
    int d = scanf("%d%d%d", &a, &b, &c);
\end{lstlisting}

其中 \lstinline{0}、\lstinline{a++}、\lstinline{++a} 都是表达式。一个表达式包含以下要素：

\begin{itemize}
    \item 对现有变量、输入输出进行的操作（也可以不做操作）；
    \item 返回值（可以为空）。
\end{itemize}

以 \lstinline{a++} 举例，它进行的操作就是将变量 \lstinline{a} 的值加一，返回值是加一前的变量的值。

\lstinline{scanf("%d%d%d", &a, &b, &c)} 也是表达式，
它的操作是读取三个变量，返回值是读取到的变量的数量。

表达式可以通过运算符、函数等组合成新的表达式：

\begin{lstlisting}[language=C]
    int e = (scanf("%d%d%d", &a, &b, &c) + ++a * b--) * c++ / --d ^ abs(a + b);
\end{lstlisting}

\subsection{函数的定义和调用}

函数的定义包括以下要素：

\begin{itemize}
    \item 参数列表；
    \item 函数体，即进行的操作；
    \item 返回值类型。
\end{itemize}

可以看出函数和表达式是很像的。

调用时需要传递参数，有返回值的话也可以使用返回值做进一步处理。

举个例子：

\label{code:2.func.define_and_invoke}

\begin{lstlisting}[language=C]
    #include <stdio.h>

    int global_arr[10], u = 0;

    float func(int x, float y, int arr[]) {
        for (int i = 0; i < 10; i++) {
            scanf("%d", &arr[i]); // 函数中不仅可以做运算，还可以输入输出
        }
        float res = ((float)x) / 10 + x / 7 + y; 
        for (int i = 0; i < 10; i++) {
            res += ++arr[i]; // 这里的操作会影响到数组中的元素
        }
        x++; // 操作后的 x 会变化，但不会改变 main 中 a 的值
        u++;
        printf("x, u after change: %d, %d\n", x, u);
        return res; // 函数返回类型不是 void 的话，就必须有返回值
    }

    int main() {
        int a;
        float b;
        scanf("%d%f", &a, &b);
        float result = func(a, b, global_arr) + a; // 调用，并接收返回值做进一步运算
        printf("Result: %.2f\n", result);
        for (int i = 0; i < 10; i++) {
            printf("global_arr[%d] = %d\n", i, global_arr[i]); // 输出会显示数组里的元素都加了 1
        }
        print("a after func: %d\n", a); // a 不会因为 x++ 而变化
        print("u after func: %d\n", u); // 但 u 会变化
        return 0;
    }
\end{lstlisting}

\begin{itemize}
    \item 传递的参数在函数内部和外部被看作是不同的变量，
    内部改变并不会影响到函数外，但对全局变量的操作是有效的；
    \item 数组也可以作为参数传递，但和单一变量参数不同，函数中数组内元素的改变会影响到函数外；
    可以这么理解：传递的数组包含的是数组的位置信息，而不是数组内部元素的值。
\end{itemize}

\section{课程前瞻}

目前我是这么想的，现在基础知识差不多了，对于写大多数简单的题目已经没问题了。
接下来一段时间先带你刷刷题，把基础打好。
我们从这个题单开始：
\href{https://github.com/ankitvashisht12/100-questions-of-codeforces}{CF 新手 100 题}。
每周你就在这个题单里面随便找几个题做一下，具体多少数量我现在也不清楚，因为刷题期间实际讲课节奏我心里也没谱，
等刷起来以后我们才知道。

要刷到什么程度？大概就是你遇到这个难度的题都能做出来就可以了。

由于前面我只关注你能快速上手实践，很多基础内容都忽略了，刷题期间我会顺带把这部分补上。
这些内容也很重要，之所以前面不讲是因为这些你在大多数教材上都能看到，
而且对于“能快速上手实践”其实关系不是很大。
然后还会讲一些绝大多数国内外课程都不会教的东西。

刷题阶段完成后，接下来就开始一些简单的算法和数据结构，你不去打 ICPC 之类的竞赛的话，不用学的那么变态，掌握基础就行了。

\begin{problemset}[练习]
    \item 为 程序 \ref{code:2.func.define_and_invoke} 随机构造几组输入，
    先手动计算输出，然后运行程序看一下和你的预期是否一致；如果不一致，看看你哪里理解错了。
\end{problemset}

\end{document}
